\lysh{16808_SOLUTION}{1}
\fontsize{12pt}{14pt}\selectfont
\noindent ΛΥΣΗ

\noindent α) Το μέσο $K$ του τμήματος $AB$ έχει συντεταγμένες $\left(\frac{x_A+x_B}{2}, \frac{y_A+y_B}{2}\right)$.
Άρα $K\left(\frac{-8+4}{2}, \frac{1+5}{2}\right) = (-2, 3)$.

\noindent β) Αρκεί να δείξουμε ότι, το μέσο $K$ του τμήματος $AB$ απέχει από το σημείο $\Gamma$ απόσταση ίση με το μισό του τμήματος $AB$. Δηλαδή $K\Gamma = \frac{AB}{2}$.

\noindent Το μήκος του τμήματος $AB$ είναι
\[ (AB) = \sqrt{(-8-4)^2 + (1-5)^2} = \sqrt{144 + 16} = \sqrt{160} = 4\sqrt{10} . \]

\noindent Το μήκος του τμήματος $K\Gamma$ είναι
\[ (K\Gamma) = \sqrt{(-2+4)^2 + (3-9)^2} = \sqrt{4 + 36} = \sqrt{40} = 2\sqrt{10} = \frac{AB}{2} . \]

\noindent γ) Το κέντρο του κύκλου είναι το σημείο $K(-2,3)$ και η ακτίνα του είναι $\rho = \frac{AB}{2} = 2\sqrt{10}$.

\noindent Άρα η εξίσωση του κύκλου είναι
\[ (C): (x+2)^2 + (y-3)^2 = 40. \]
\endlysh{}